\documentclass[a4paper,11pt]{article}
\usepackage[english]{babel}
\usepackage[utf8]{inputenc}
\usepackage[T1]{fontenc}
\usepackage[a4paper,margin=1.5cm,top=1.2cm,bottom=1.2cm]{geometry}
\usepackage{graphicx}
\usepackage{amsmath}
\usepackage{booktabs}
\usepackage{siunitx}
\usepackage{float}
\usepackage{microtype}
\usepackage{newtxtext,newtxmath}
\usepackage{setspace}
\usepackage{fancyhdr}
\setstretch{1.15}

\bibliographystyle{apsrev}
\setcounter{secnumdepth}{3}
\pagestyle{fancy}
\fancyhf{}
\fancyfoot[C]{\thepage}
\renewcommand{\headrulewidth}{0pt}
\renewcommand{\footrulewidth}{0pt}
\setlength{\headheight}{14pt}
\setlength{\footskip}{14pt}

\begin{document}

\title{Determinacion de la longitud de onda mediante un biprisma de Fresnel}
\author{Practicas de Optica}
\maketitle

\begin{abstract}
Se determina experimentalmente la longitud de onda de la luz mediante un biprisma
de Fresnel. El analisis utiliza las distancias medidas entre las fuentes virtuales
y la pantalla, junto con la separacion entre franjas. Se propagan las incertidumbres
tipo A y tipo B mediante la biblioteca \texttt{uncertainties}. Los resultados se
comparan con el valor teorico de $589\,\mathrm{nm}$.
\end{abstract}

\section{Introduccion y objetivos}

\subsection{Fundamento teorico}
El biprisma produce dos fuentes virtuales coherentes separadas una distancia $d$.
Si $D$ es la distancia desde las fuentes virtuales hasta la pantalla y $s$ es la
separacion entre franjas consecutivas, la longitud de onda se obtiene mediante
\begin{equation}
\lambda = \frac{s d}{D}.
\end{equation}
Para la geometria empleada, la distancia entre las fuentes virtuales se calcula
como
\begin{equation}
d = \sqrt{d'_1 d'_2},
\end{equation}
con $d'_1$ y $d'_2$ las distancias medidas desde la lente a las imagenes virtuales.
La distancia $D$ se obtiene de
\begin{equation}
D = f'\left(2 - \frac{d^2 + {d'_1}^2}{d d'_1}\right).
\end{equation}

% Example equations block (replace with experiment-specific model)
\subsection{Metodologia de medida}
Se midieron tres valores de $s$, $d'_1$ y $d'_2$ en cada parte del experimento.
En las partes 1 y 2 se uso una distancia focal $f'=0.15\,\mathrm{m}$, mientras
que en la parte 3 se uso $f'=0.20\,\mathrm{m}$. En las partes 2 y 3, $s$ representa
una interfranja; en la parte 1 se midieron tres interfranjas, por lo que se uso
$\lambda=s d/(3D)$.

\section{Materiales y metodos}

\subsection{Montaje experimental}
El montaje esta formado por una fuente de luz de sodio, un biprisma de Fresnel,
una lente convergente y una pantalla de observacion. Las posiciones se registraron
en unidades de $10^{-6}\,\mathrm{m}$, tal como se introdujeron en el script de analisis.

\subsection{Adquisicion y analisis de datos}
Los datos se introdujeron en los arrays de \texttt{p4.py}. La incertidumbre de
cada magnitud se calculo combinando en cuadratura una contribucion instrumental
tipo B y la incertidumbre estadistica tipo A, $u_A=s/\sqrt{n}$, usando la desviacion
estandar muestral ($\mathrm{ddof}=1$). Las incertidumbres se propagaron con
\texttt{unumpy}.

\section{Resultados y analisis de datos}

\subsection{Resultados de la longitud de onda}
La Tabla~\ref{tab:lambda} resume los resultados obtenidos para cada medida.

\begin{table}[H]
\centering
\begin{tabular}{lccc}
\toprule
\textbf{Parte} & $\lambda_1$ & $\lambda_2$ & $\lambda_3$ \\
\midrule
Parte 1 & $577\pm87$ & $594\pm95$ & $569\pm92$ \\
Parte 2 & $693\pm33$ & $808\pm33$ & $750\pm33$ \\
Parte 3 & $491\pm130$ & $508\pm122$ & $505\pm146$ \\
\bottomrule
\end{tabular}
\caption{Longitudes de onda individuales en nm.}
\label{tab:lambda}
\end{table}

\subsection{Promedios}
\begin{table}[H]
\centering
\begin{tabular}{lc}
\toprule
\textbf{Parte} & \textbf{Promedio de $\lambda$ (nm)} \\
\midrule
Parte 1 & $580\pm53$ \\
Parte 2 & $750\pm39$ \\
Parte 3 & $501\pm81$ \\
\bottomrule
\end{tabular}
\caption{Promedios calculados combinando la dispersion entre medidas y sus incertidumbres propagadas.}
\label{tab:promedios}
\end{table}

El valor teorico de referencia es $\lambda_\mathrm{teo}=589\,\mathrm{nm}$.
El resultado de la primera parte, $580\pm53\,\mathrm{nm}$, es compatible con
este valor dentro de la incertidumbre experimental.

\section{Discusion y conclusiones}
La primera parte presenta el mejor acuerdo con el valor teorico de la fuente de
sodio. Las diferencias entre las partes 2 y 3 pueden deberse a la sensibilidad
de $D$ frente a las distancias medidas y a la incertidumbre relativamente grande
de la mirilla usada en la tercera parte. La propagacion con \texttt{unumpy} incluye
las incertidumbres de $s$, $d'_1$ y $d'_2$ en $d$, $D$ y $\lambda$.

\begin{equation}
\lambda_1 = (580\pm53)\,\mathrm{nm}, \qquad
\lambda_2 = (750\pm39)\,\mathrm{nm}, \qquad
\lambda_3 = (501\pm81)\,\mathrm{nm}.
\end{equation}

\end{document}
